%% ma: L11-B16-0001
%% lop: 11
%% chuong: 5
%% bai: 16
%% dang: 11.16.1
%% loai: TL
%% muc: TH
%% dap_an: "-"
%% nguon: "tai-lieu-hoc-tap-mon-toan-lop-11-theo-chuong-trinh-giao-duc-pho-thong-2018.pdf (Lớp Toán Thầy Duy, Phạm Lê Duy, THPT Chu Văn An), Chương 3 Bài 2, Dạng 1, Ví dụ 1 (tr.351)"
%% kien_thuc: "Lớp 11 Bài 16 (giới hạn của hàm số)"
%% trang_thai: chinh-thuc
%% ghi_chu: "Đã duyệt 03/11/2026."
\begin{ex}
Tìm giới hạn

a) $\lim\limits_{x\to4}\dfrac{x^2+2x-8}{x^2+4x}$;

b) $\lim\limits_{x\to\frac12}\dfrac{2x^2-5x+2}{1-2x}$;

c) $\lim\limits_{x\to-2}\dfrac{2x^2-5x+2}{x^2+x-6}$;

d) $\lim\limits_{x\to-1}\dfrac{1+x^3}{1-x^2}$.
\loigiai{\sloppy a) Tại $x=4$ mẫu $=32\ne0$ nên giới hạn $=\dfrac{16+8-8}{16+16}=\dfrac12$.

b) $\dfrac{2x^2-5x+2}{1-2x}=\dfrac{(2x-1)(x-2)}{-(2x-1)}=2-x\to\dfrac32$.

c) Tại $x=-2$: tử $=20$, mẫu $=-4\ne0$ nên giới hạn $=-5$.

d) $\dfrac{1+x^3}{1-x^2}=\dfrac{(1+x)(1-x+x^2)}{(1-x)(1+x)}=\dfrac{1-x+x^2}{1-x}\to\dfrac32$.}
\end{ex}
